\clearpage
\VerdierPhysicalPageAnchor{0206}%
\begingroup\small\setlength{\parskip}{0pt}\setlength{\emergencystretch}{1.5em}
\setlength{\abovedisplayskip}{4pt plus 2pt minus 2pt}%
\setlength{\belowdisplayskip}{4pt plus 2pt minus 2pt}%

\VerdierSemanticAnchor{ch3d-2-4-5}
\noindent\textbf{注 2.4.5.}\quad
a) 设 $\mathcal A$ 是阿贝尔范畴，其每个对象都同构于投射对象的商对象。
投射型复形不一定是 $Q$-左自由的。

\VerdierSemanticAnchor{ch3d-2-4-5-example}
\emph{例}：设 $\mathcal A$ 是 $\mathbb Z/4\mathbb Z$-模范畴。复形
\[
\cdots\longrightarrow\mathbb Z/4\mathbb Z
\longrightarrow\mathbb Z/4\mathbb Z
\longrightarrow\mathbb Z/4\mathbb Z
\longrightarrow\mathbb Z/4\mathbb Z
\longrightarrow\mathbb Z/4\mathbb Z\longrightarrow\cdots
\]
的所有微分均为乘以 $2$；它无环且为投射型复形，但并不同伦于零。同一个复形
也给出了一个不同伦于零的无环内射型复形的例子。

\VerdierSemanticAnchor{ch3d-2-4-5-b}
b) 即使 $\mathcal A$ 的对象的投射维数无上界，也可能存在 $Q$-左自由的
投射型复形，其在 $D(\mathcal A)$ 中的像不属于 $D^{-}(\mathcal A)$。
例如，Cartan--Eilenberg 投射复形是 $Q$-左自由对象。更一般地，
$Q$-左自由对象的任意直和仍是 $Q$-左自由的。

\VerdierSemanticAnchor{ch3d-2-4-6}
\medskip\noindent\textbf{2.4.6.}\quad
把本节结果应用于范畴 $\mathcal A$ 的反范畴所得的陈述，留给读者表述。
\endgroup

